خیر، لزوماً این‌طور نیست.

متخصص شدن در یک حوزه‌ی خاص از مهندسی نرم‌افزار (مثلاً فقط پایگاه‌داده، فرانت‌اند یا الگوریتم‌ها) به این معنا نیست که آن فرد حتماً معمار نرم‌افزار خوبی هم خواهد شد. معماری نرم‌افزار نیازمند دیدی گسترده‌تر از صرفاً تسلط بر یک موضوع است.

یک معمار خوب باید بتواند بین بخش‌های مختلف سیستم ارتباط برقرار کند، تصمیم‌های کلان بگیرد، و trade-off ها (مثل کارایی، مقیاس‌پذیری، نگهداری‌پذیری و امنیت) را به‌درستی مدیریت کند. این مهارت‌ها معمولاً از تجربه در حوزه‌های مختلف، درک کلی از سیستم‌ها، و توانایی تحلیل مسائل در سطح بالا به دست می‌آیند، نه فقط از عمیق شدن در یک زمینه‌ی خاص.

بنابراین، تخصص عمیق در یک حوزه می‌تواند یک مزیت باشد، اما برای تبدیل شدن به یک معمار خوب کافی نیست؛ بلکه باید در کنار آن، دید سیستمی و تجربه‌ی گسترده نیز وجود داشته باشد.